% ============================================================
% SECTION 7b — Spatial Diagnostics Subsection
% Home: SIGSPATIAL (canonical spatial-statistics vocabulary).
% NeurIPS dataset paper references this, does not duplicate figures.
%
% Four diagnostics, one narrative arc:
%   1. Errors are spatially structured (global Moran + Geary)
%   2. LISA localizes WHERE the structure lives (Fig 5)
%   3. The relationship is non-stationary, explaining WHY (GWR, Fig 6)
%   4. The headline survives alternative spatial blocking (regionalization)
%
% ALL NUMBERS ARE \PH{PLACEHOLDERS} UNTIL DATA LOCK.
% ============================================================

\newcommand{\PH}[1]{\textcolor{red}{[#1]}}  % Redundant if already defined

\subsection{Spatial Validity Diagnostics}
\label{sec:spatial-diagnostics}

We report four spatial diagnostics that collectively characterize
geometric overfitting and the effectiveness of spatially-blocked
validation. Each answers a distinct question about spatial structure
in the residuals; together they form the ``geometric overfitting
$\to$ spatially-blocked validation'' thread
(Table~\ref{tab:residual-diagnostics}, Figures~\ref{fig:lisa-triptych}
and~\ref{fig:gwr-local-r2}).


% ===================================================================
% TABLE 1: Residual-Diagnostics Table
% One row per (scenario, target, level). Primary cell in main text;
% full grid in Appendix Table~\ref{app:tab:residual-all}.
% ===================================================================

\begin{table}[t]
\centering
\caption{Residual spatial-autocorrelation diagnostics, primary cell
  (\texttt{houston}, \texttt{obs\_nfip\_event\_claims}).
  Moran's $I > 0$ with $p < 0.05$ indicates residuals are
  spatially clustered; Geary's $C < 1$ corroborates at short range.
  Cluster fractions are BH-FDR corrected ($q = 0.05$).
  All numbers are \PH{placeholders} pending data lock.}
\label{tab:residual-diagnostics}
\small
\begin{tabular}{@{}lccccccc@{}}
\toprule
\textbf{Level} &
  \textbf{Moran's $I$} &
  \textbf{Geary's $C$} &
  \textbf{$p_C$} &
  \textbf{\% HH} &
  \textbf{\% LL} &
  \textbf{\% outlier} &
  \textbf{\% sig.} \\
\midrule
R0 (static)       & \PH{0.xx} & \PH{0.xx} & \PH{0.xx} & \PH{xx} & \PH{xx} & \PH{xx} & \PH{xx} \\
R1 (+ hydrology)  & \PH{0.xx} & \PH{0.xx} & \PH{0.xx} & \PH{xx} & \PH{xx} & \PH{xx} & \PH{xx} \\
R2 (+ temporal)   & \PH{0.xx} & \PH{0.xx} & \PH{0.xx} & \PH{xx} & \PH{xx} & \PH{xx} & \PH{xx} \\
\bottomrule
\end{tabular}
\end{table}


% ===================================================================
% FIGURE 5: LISA Residual-Cluster Triptych
% Three panels: R0 -> R1 -> R2.
% Hedges baked into caption so epistemics precede numbers.
% ===================================================================

\begin{figure*}[t]
\centering
\includegraphics[width=\textwidth]{figures/fig5_lisa_triptych_houston.pdf}
\caption{%
  Local Moran's $I$ residual clusters for the primary cell
  (\texttt{houston}, \texttt{obs\_nfip\_event\_claims}, \texttt{histgbdt}),
  three representation levels.
  ZCTAs colored by LISA cluster label (BH-FDR $q = 0.05$, 999 permutations):
  HH~(hot spot), LL~(cold spot), HL/LH~(spatial outliers), ns~(not significant).
  %
  Global Moran's $I$ and significant-cluster fraction annotated per panel.
  %
  \textbf{Finding:} Local Moran's $I$ on R0 residuals identifies
  significant positive-autocorrelation clusters in \PH{N}\% of
  primary-cell ZCTAs; under R1's spatial features these attenuate
  to \PH{M}\%.
  %
  \textbf{Scope:} Cross-cell attenuation is descriptive---clusters
  attenuated in \PH{k} of 9 cells (Clopper--Pearson 95\% CI
  [\PH{lo}, \PH{hi}], reported descriptively; $n = 9$ precludes
  a formal test).
  Attenuation is consistent with the W-matrix features absorbing
  the spatial structure; it is not causal proof.
  The spatially-blocked design is causal at the pipeline level, not
  the individual-feature level.%
}
\label{fig:lisa-triptych}
\end{figure*}


% ===================================================================
% FIGURE 6: GWR Local-R2 Choropleth
% Structural description only. Not predictive performance.
% ===================================================================

\begin{figure}[t]
\centering
\includegraphics[width=\columnwidth]{figures/fig6_gwr_local_r2_houston.pdf}
\caption{%
  Geographically Weighted Regression local $R^2$ for the primary cell
  (\texttt{houston}, \texttt{obs\_nfip\_event\_claims}).
  Left: per-ZCTA local $R^2$ from GWR fit (EPSG:5070, adaptive bisquare
  kernel, bandwidth selected by \texttt{AICc}).
  Right: per-feature coefficient spread (std$/|$mean$|$); features above
  the dashed threshold ($= 0.5$) exhibit non-stationary relationships.
  %
  \textbf{Structural description:} GWR improves on global OLS
  ($\Delta$AICc $=$ \PH{xx}), and local $R^2$ ranges from
  \PH{lo} to \PH{hi}, indicating the feature$\to$claims
  relationship is non-stationary across space.
  This is a description of the full cross-section using all data;
  GWR makes no generalization claim.
  %
  \textbf{Triangulation} (observed association, not tested, not in
  the Holm--Bonferroni family): low-local-$R^2$ regions coincide
  with LISA hot-spot clusters and with the highest R0$\to$R1 uplift
  (rank agreement; Spearman $\rho =$ \PH{xx}, bootstrap 95\% CI
  [\PH{lo}, \PH{hi}]).
  At $n \approx$ \PH{4--5} scenarios the CI is near-useless and we
  note it as suggestive only.%
}
\label{fig:gwr-local-r2}
\end{figure}


% ===================================================================
% Geary's C -- one sentence, corroborating
% ===================================================================

\paragraph{Short-range dissimilarity.}
Geary's $C$ corroborates the Moran's $I$ finding:
$C =$ \PH{0.xx} at R0 ($p =$ \PH{0.xx}), consistent in direction
($C < 1$ indicates positive spatial autocorrelation).
See Table~\ref{tab:residual-diagnostics} for all levels.


% ===================================================================
% TABLE 2: Robustness Table (Appendix R)
% One row per compared cell. Header carries overlap counts.
% ===================================================================

\begin{table}[t]
\centering
\caption{Regionalization robustness: R0$\to$R1 uplift direction under
  county-blocked versus region-blocked (Skater, $K = 5$) folds.
  Compared \PH{m} of \PH{n} cells (county: \PH{n\_county},
  region: \PH{n\_region}; \PH{dropped keys if any}).
  The county-blocked Wilcoxon remains the registered primary inference;
  this is a sensitivity check on block geometry.
  %
  \textbf{wlag NaN confound:} \texttt{wlag\_nfip\_claims} was absent for
  \PH{X}\% of region-fold test ZCTAs versus \PH{Y}\% under county
  folds, a consequence of contiguous holdout (interior test ZCTAs lose
  all training neighbors).
  Any uplift-direction difference may partly trace to differential
  feature availability rather than pure block geometry.%
}
\label{tab:robustness}
\small
\begin{tabular}{@{}llccccl@{}}
\toprule
\textbf{Scenario} & \textbf{Target} &
  \textbf{Uplift (county)} & \textbf{$p_W$ (county)} &
  \textbf{Uplift (region)} & \textbf{$p_W$ (region)} &
  \textbf{Agrees?} \\
\midrule
\PH{houston} & \PH{claims} & \PH{+0.xx} & \PH{0.xx} & \PH{+0.xx} & \PH{0.xx} & \PH{yes} \\
\PH{houston} & \PH{311}    & \PH{+0.xx} & \PH{0.xx} & \PH{+0.xx} & \PH{0.xx} & \PH{yes} \\
\PH{houston} & \PH{hwm}    & \PH{+0.xx} & \PH{0.xx} & \PH{+0.xx} & \PH{0.xx} & \PH{--}  \\
\PH{\ldots}  & \PH{\ldots} & \PH{\ldots} & \PH{\ldots} & \PH{\ldots} & \PH{\ldots} & \PH{\ldots} \\
\midrule
\multicolumn{6}{@{}l}{\textbf{Direction agreement:}} & \PH{k / m} \\
\bottomrule
\end{tabular}
\end{table}


% ===================================================================
% Regionalization paragraph
% ===================================================================

\paragraph{Robustness to block geometry (Appendix~R).}
Under an alternative hydrologically-coherent block definition
(Skater regionalization on structural geography%
---\texttt{flood\_pct\_zone\_a}, \texttt{twi\_acc\_twi},
\texttt{slope\_basin\_slope}%
---$K = 5$), the sign of R0$\to$R1 uplift agrees with county blocking
in \PH{k of m} compared cells (Table~\ref{tab:robustness}).
Where disagreement occurs, we report it as sensitivity to block geometry,
not as failure of either blocking scheme.

\textbf{Partial overlap:}
The comparison covers \PH{m} of 9 cells;
\PH{dropped cells and reason, e.g.\ ``SW Florida dropped due to
Skater degeneracy (singleton region)''}.
The asymmetry is visible in the artifact (\texttt{n\_county\_cells},
\texttt{n\_region\_cells}, \texttt{dropped\_*} fields in the comparison
JSON) and does not inflate the reported agreement fraction.


% ===================================================================
% Narrative arc (closing paragraph)
% ===================================================================

\paragraph{Summary.}
Residuals are spatially structured at R0 (Moran's $I$, Geary's $C$);
LISA localizes \emph{where} the structure concentrates
(Figure~\ref{fig:lisa-triptych});
the underlying relationship is non-stationary, explaining \emph{why}
spatial features help (GWR, Figure~\ref{fig:gwr-local-r2});
R1's W-matrix features attenuate the clusters;
and the headline uplift survives an alternative spatial blocking
(Table~\ref{tab:robustness}).
This sequence instantiates the ``geometric overfitting $\to$
spatially-blocked validation'' thread that the benchmark is designed
to diagnose.
